% This is text associated with the open textbook "Open Electromagnetics"
% See immediately below for author, date
% This work is licensed under the Creative Commons Attribution-ShareAlike 4.0 International License. (CC BY-SA 4.0) 
% To view a copy of the license, visit http://creativecommons.org/licenses/by-sa/4.0/

%ORB TITLE Applications of Open- and Short-Circuited Transmission Line Stubs
%ORB AUTHOR 0        # S.W. Ellingson, Virginia Tech (ellingson@vt.edu)
%ORB DATE 20180802   # yyyymmdd
%ORB ABSTRACT Applications of Open- and Short-Circuited Transmission Lines

\label{m0145_Applications_of_Open_and_Short-Circuited_Transmission_Lines} % <-- Must match name of this file and module directory name.  Don't mess with this!
\index{stub!applications}

The theory of open- and short-circuited transmission lines -- often referred to as \emph{stubs} -- was addressed in Section~\ref{m0088_Input_Impedance_for_Open_and_Short_Circuit_Terminations}. 
These structures have important and wide-ranging applications. 

In particular, these structures can be used to replace discrete inductors and capacitors in certain applications.  
To see this, consider the short-circuited line (Figure~\ref{m0088_fZin}(b) of Section~\ref{m0088_Input_Impedance_for_Open_and_Short_Circuit_Terminations}). Note that each value of $l$ that is less than $\lambda/4$ corresponds to a particular positive reactance; i.e., the transmission line ``looks'' like an inductor.
Also note that lengths between $\lambda/4$ and $\lambda/2$ result in reactances that are negative; i.e., the transmission line ``looks'' like a capacitor.
Thus, it is possible to replace an inductor or capacitor with a short-circuited transmission line of the appropriate length.
The input impedance of such a transmission line is identical to that of the inductor or capacitor at the design frequency. 
The variation of reactance with respect to frequency will \emph{not} be identical, which may or may not be a concern depending on the bandwidth and frequency response requirements of the application.
Open-circuited lines may be used in a similar way. 

This property of open- and short-circuited transmission lines makes it possible to implement impedance matching 
\index{impedance!matching}
circuits (see Section~\ref{m0094_Single-Stub_Matching}), filters, 
\index{filter}
and other devices entirely from transmission lines, with fewer or no discrete inductors or capacitors required.
Transmission lines do not suffer the performance limitations of discrete devices at high frequencies and are less expensive.
A drawback of transmission line stubs in this application is that the lines are typically much larger than the discrete devices they are intended to replace.

%ORB EXAMPLE
~\\ \begin{example} 
\begin{mdframed}[backgroundcolor=brown!20,linecolor=brown!20]\raggedright
Emitter Induction 
\index{emitter induction}
Using Short-Circuited Line. 
\\~\\
In the design of low-noise amplifiers using bipolar transistors in common-emitter configuration, it is often useful to introduce a little inductance between the emitter and ground.
This is known as ``inductive degeneration,'' ``emitter induction,'' or sometimes by other names.
It can be difficult to find suitable inductors, especially for operation in the UHF band and higher.
However, a microstrip 
\index{microstrip}
line can be used to achieve the desired inductive impedance.
Determine the length of a stub that implements a 2.2~nH inductance at 6~GHz using microstrip line with characteristic impedance $50~\Omega$ and phase velocity $0.6c$. 

{\bf Solution.}
At the design frequency, the impedance looking into this section of line from the emitter should be equal to that of a 2.2~nH inductor, which is 
$+j\omega L = +j2\pi fL = +j82.9~\Omega$.
The input impedance of a short-circuited stub of length $l$ which is grounded (thus, short-circuited) at the opposite end is $+jZ_0\tan\beta l$ (Section~\ref{m0088_Input_Impedance_for_Open_and_Short_Circuit_Terminations}).  
Setting this equal to $+j82.9~\Omega$ and noting that $Z_0=50~\Omega$, we find that $\beta l \cong 1.028$~rad.
The phase propagation constant is (Section~\ref{m0080_Wave_Propagation_on_a_Transmission_Line}):
%
\begin{equation}
\beta = \frac{\omega}{v_p} = \frac{2\pi f}{0.6 c} \cong 209.4~\mbox{rad/m}
\end{equation}
%
Therefore, the length of the microstrip line is $l = \left(\beta l\right) / \beta \cong 4.9$~mm.

\end{mdframed}\end{example}

%%%%%%%%%%%%%%%%%%%%%%%%%%%%%%%%

{\bf Additional Reading:}
\begin{itemize}
\item ``\href{https://en.wikipedia.org/wiki/Stub_(electronics)}{Stub (electronics)}'' on Wikipedia.
\end{itemize}






